% preambolo.tex — del gruppo: si modifica a mano. Gli script non lo toccano.
\usepackage[T1]{fontenc}
\usepackage[utf8]{inputenc}
\usepackage{lmodern}
\usepackage[english]{babel}
\usepackage{csquotes}
\usepackage[a4paper,margin=2.0cm]{geometry}
\usepackage{setspace}\setstretch{1.05}
\usepackage{microtype}
\usepackage{newunicodechar}         % simboli Unicode nel testo e nelle didascalie
\newunicodechar{β}{\ensuremath{\beta}}\newunicodechar{ρ}{\ensuremath{\rho}}
\newunicodechar{θ}{\ensuremath{\theta}}\newunicodechar{χ}{\ensuremath{\chi}}
\newunicodechar{×}{\ensuremath{\times}}\newunicodechar{−}{\ensuremath{-}}
\newunicodechar{≈}{\ensuremath{\approx}}\newunicodechar{≤}{\ensuremath{\leq}}
\newunicodechar{≥}{\ensuremath{\geq}}\newunicodechar{→}{\ensuremath{\rightarrow}}
\newunicodechar{±}{\ensuremath{\pm}}\newunicodechar{′}{\ensuremath{'}}
\newunicodechar{²}{\ensuremath{^2}}\newunicodechar{Σ}{\ensuremath{\Sigma}}\newunicodechar{≠}{\ensuremath{\neq}}\newunicodechar{–}{--}
\usepackage{amsmath,amssymb}
\usepackage{graphicx}
\usepackage{tikz}\usetikzlibrary{decorations.pathreplacing}
\usepackage{booktabs}
\usepackage{array}
\usepackage{caption}\captionsetup{font=small,labelfont=bf}
\usepackage[section]{placeins}     % i floats non scavalcano la sezione successiva
\usepackage{flafter}               % ...e non salgono sopra il punto in cui sono richiamati
\usepackage{xcolor}
\input{gen/bib}                     % biblatex-apa se c'e', altrimenti natbib (scritto da assemble.py)
\usepackage[hidelinks]{hyperref}
\graphicspath{{figures/}}
\setlength{\parskip}{0.3em}
% \buco{...}: segnaposto di un buco dichiarato (solo nelle bozze: il PDF finale non ne ha)
\newcommand{\buco}[1]{\fcolorbox{red}{red!8}{\parbox{0.9\linewidth}{\small\textbf{BUCO} #1}}}
% --- Floats fra i paragrafi (Lele, 23/11): nessuna figura o tabella spezza un paragrafo.
% Ogni float del corpo e' composto in un box e messo in coda; dopo ogni paragrafo
% \tentafloats lo inserisce se ci sta nello spazio rimasto nella pagina, altrimenti il testo
% continua e il float aspetta la fine del paragrafo successivo. \svuotafloats li mette tutti
% (a inizio sezione e a fine corpo).
\makeatletter
\newcount\fq@coda \newcount\fq@testa
\newdimen\fq@spazio
\newenvironment{codafloat}[1]{%
  \global\advance\fq@coda\@ne
  \expandafter\newbox\csname fq@box\the\fq@coda\endcsname
  \global\setbox\csname fq@box\the\fq@coda\endcsname\vbox\bgroup
    \hsize\linewidth \@parboxrestore \captionsetup{type=#1}\centering}%
  {\egroup}
\def\fq@metti{\global\advance\fq@testa\@ne
  \par\addvspace{\intextsep}%
  \nointerlineskip\box\csname fq@box\the\fq@testa\endcsname
  \par\addvspace{\intextsep}}
\def\fq@altezza{\dimexpr\ht\csname fq@box\the\numexpr\fq@testa+1\endcsname
  +\dp\csname fq@box\the\numexpr\fq@testa+1\endcsname+2\intextsep\relax}
\newif\iffq@ancora
\def\tentafloats{\par\fq@ancoratrue
  \loop
    \ifnum\fq@testa<\fq@coda
      \ifdim\pagegoal=\maxdimen \fq@metti
      \else
        \fq@spazio=\dimexpr\pagegoal-\pagetotal\relax
        \ifdim\fq@altezza<\fq@spazio \fq@metti\else\fq@ancorafalse\fi
      \fi
    \else\fq@ancorafalse\fi
  \iffq@ancora\repeat}
\def\svuotafloats{\par
  \ifnum\fq@testa<\fq@coda \fq@metti\expandafter\svuotafloats\fi}
\makeatother
